% ============================================================================
% Chapter 8: L₀/δ Scale Ratio
% ============================================================================
% Source: DERIVE_L0_DELTA_PI_SQUARED_V2.md
% Status: FILLED from SoT
% Contamination check: PASSED (no SM terms)
% ============================================================================

\chapter{$\Lzero/\bdelta$ Scale Ratio}
\label{ch:L0delta}

\source{DERIVE\_L0\_DELTA\_PI\_SQUARED\_V2.md}

\begin{abstract}
Chapter~\ref{ch:kappa} established $\kappa = 2\pi$ from homotopy. The Euclidean
action $\SE/\hbar = \kappa \times (\Lzero/\bdelta)$ still requires the geometric
ratio $\Lzero/\bdelta$. This chapter examines the hypothesis that
$\Lzero/\bdelta \approx \pi^2$. We present geometric motivations from standing
wave quantization and resonance conditions, but the result remains [P]---a
physically motivated conjecture, not a rigorous derivation from the 5D action.
\end{abstract}

% ----------------------------------------------------------------------------
% CHAPTER SPINE
% ----------------------------------------------------------------------------
\paragraph{Chapter Spine.}
\textbf{Purpose:} Establish the geometric ratio $L_0/\delta \approx \pi^2$
that completes the Euclidean action formula.
\textbf{Inputs:} (i)~$\SE$ structure from Ch.~\ref{ch:instanton}~[Dc],
(ii)~$\kappa = 2\pi$ from Ch.~\ref{ch:kappa}~[Dc], (iii)~$\delta = \hbar/(2m_p c)$~[I].
\textbf{Outputs:} (i)~$L_0/\delta = \pi^2$~[P] (hypothesis, not derivation),
(ii)~$\SE/\hbar = 2\pi^3 \approx 62$.
\textbf{Dependencies:} Ch.~\ref{ch:kappa} ($\kappa = 2\pi$).
\textbf{Forward links:} $L_0/\delta$ to Ch.~\ref{ch:tau_n} ($\taun$ assembly).
\textbf{Epistemic note:} This ratio is \emph{postulated}~[P], not derived~[Der].
A first-principles derivation remains OPEN.

% ============================================================================
\section{Two Length Scales}
\label{sec:L0delta:scales}

\subsection{The Junction Extent $\Lzero$}

The quantity $\Lzero$ is the \emph{characteristic extent} of the junction
configuration in the fifth dimension:

\begin{definition}[Junction Extent $\Lzero$]
    $\Lzero$ is the length scale over which the Y-junction structure
    (anchor/metastable configuration) extends into the compact fifth dimension.
    It sets the ``size'' of the topological defect.
    \tagDc
\end{definition}

\textbf{Physical interpretation:}
\begin{itemize}
    \item $\Lzero$ is the radial extent from junction center to effective boundary
    \item For anchor: $\Lzero \approx r_p + \bdelta \approx 1\fm$ (order of magnitude)
    \item $\Lzero$ controls the energy stored in the junction configuration
\end{itemize}

\subsection{The Brane Scale $\bdelta$}

The quantity $\bdelta$ is the \emph{fundamental thickness} of the 5D brane:

\begin{definition}[Brane Thickness $\bdelta$]
    $\bdelta$ is the characteristic scale of the compact fifth dimension,
    related to the regularization of the brane worldvolume. In natural units:
    \begin{equation}
        \bdelta = \frac{\hbar}{2 m_p c} \approx 0.105\fm
        \label{eq:delta_def}
        \tagI
    \end{equation}
    where $m_p$ is the anchor mass.
\end{definition}

\textbf{Origin:} $\bdelta$ arises from Book I (5D foundations) as the Compton-scale
regularization of the brane. We treat it here as an input parameter [I].

\subsection{Dimensional Consistency}

Both $\Lzero$ and $\bdelta$ have dimensions of length. Their ratio is dimensionless:
\begin{equation}
    \frac{\Lzero}{\bdelta} = \text{(pure number)}
\end{equation}

This ratio appears in the instanton action (Chapter~\ref{ch:instanton}):
\begin{equation}
    \frac{\SE}{\hbar} = \kappa \times \frac{\Lzero}{\bdelta} = 2\pi \times \frac{\Lzero}{\bdelta}
\end{equation}

The exponential sensitivity of metastable lifetime to this ratio ($\taun \propto \exp[\SE/\hbar]$)
makes its precise value critical.

% ============================================================================
\section{The $\pi^2$ Hypothesis}
\label{sec:L0delta:hypothesis}

\subsection{Statement}

\begin{derivationbox}[title=Hypothesis: $\Lzero/\bdelta = \pi^2$]
    \textbf{Claim:} The ratio of junction extent to brane thickness is:
    \begin{equation}
        \boxed{\frac{\Lzero}{\bdelta} = \pi^2 \approx 9.87}
        \label{eq:L0delta_hypothesis}
    \end{equation}
    \tagP
\end{derivationbox}

This is a \emph{hypothesis}---geometrically motivated but not rigorously derived
from the 5D action.

\subsection{Numerical Implications}

With $\bdelta = 0.105\fm$:
\begin{equation}
    \Lzero = \pi^2 \times \bdelta = 9.87 \times 0.105\fm \approx 1.04\fm
\end{equation}

This is consistent with the observed anchor charge radius $r_p \approx 0.88\fm$
(since $\Lzero = r_p + \bdelta$).

\subsection{For Metastable Lifetime}

Combined with $\kappa = 2\pi$ (Chapter~\ref{ch:kappa}):
\begin{equation}
    \frac{\SE}{\hbar} = 2\pi \times \pi^2 = 2\pi^3 \approx 62.0
\end{equation}

This gives $\exp(\SE/\hbar) \approx 8.8 \times 10^{26}$, which (with appropriate
prefactors) yields $\taun \sim 10^3\second$---the correct order of magnitude.

\begin{edcopenbox}[title=Assumption Check: $\pi^2$ Value]
    \textbf{What is assumed:}
    \begin{itemize}
        \item The ratio $\Lzero/\bdelta$ is a pure geometric constant
        \item This constant equals $\pi^2$ (not $3\pi$, $2\pi+1$, or other)
        \item The same ratio applies to both static properties (mass) and
              dynamic processes (tunneling)
    \end{itemize}

    \textbf{What would falsify this:}
    \begin{itemize}
        \item If $\taun$ calculation with $\pi^2$ mismatches by $>10\%$ after
              accounting for prefactor $A$
        \item If independent measurement of $\Lzero$ gives
              $\Lzero/\bdelta \neq \pi^2 \pm 5\%$
    \end{itemize}
    \tagP
\end{edcopenbox}

% ============================================================================
\section{Geometric Motivation}
\label{sec:L0delta:geometry}

Several geometric arguments suggest $\Lzero/\bdelta \approx \pi^2$. None is
fully rigorous, but together they provide compelling motivation.

\subsection{Approach 1: Standing Wave in Compact Dimension}

\textbf{Physical picture:} The compact fifth dimension has circumference $2\pi R_5$.
A standing wave in this dimension has wavelength $\lambda_n = 2\pi R_5 / n$.

\begin{derivationbox}[title=Standing Wave Argument]
    For the fundamental mode ($n = 1$) to fit within the junction:
    \begin{equation}
        \Lzero = \frac{\lambda_1}{2} = \pi R_5
    \end{equation}

    \textbf{Key step [P]:} Assume the compactification radius is:
    \begin{equation}
        R_5 = \pi \bdelta
    \end{equation}

    Then:
    \begin{equation}
        \Lzero = \pi \times \pi \bdelta = \pi^2 \bdelta
    \end{equation}
    \tagP
\end{derivationbox}

\textbf{Epistemic status:} The standing wave structure is [Dc] given the $\Sone$
topology. The identification $R_5 = \pi\bdelta$ is [P]---a hypothesis relating
the compactification radius to the brane thickness.

\subsection{Approach 2: Two-Fold Winding}

\textbf{Physical picture:} The junction involves two independent sources of $\pi$:
\begin{enumerate}
    \item \textbf{Radial quantization:} A half-wavelength fits in the radial
          direction, contributing factor $\pi$
    \item \textbf{Angular quantization:} A full winding around $\Sone$
          contributes factor $\pi$ (from $\oint d\theta / 2 = \pi$)
\end{enumerate}

\begin{derivationbox}[title=Two-Fold Winding Argument]
    \textbf{Radial condition:} The radial wavefunction has a node at the center
    and at the boundary:
    \begin{equation}
        k_r \Lzero = \pi \quad \Rightarrow \quad k_r = \frac{\pi}{\Lzero}
    \end{equation}

    \textbf{Angular condition:} The angular wavefunction winds once:
    \begin{equation}
        k_\theta \times 2\pi R_5 = 2\pi \quad \Rightarrow \quad k_\theta = \frac{1}{R_5}
    \end{equation}

    \textbf{Resonance [P]:} Matching $k_r = k_\theta$ gives:
    \begin{equation}
        \frac{\pi}{\Lzero} = \frac{1}{R_5} \quad \Rightarrow \quad \Lzero = \pi R_5
    \end{equation}

    With $R_5 = \pi\bdelta$:
    \begin{equation}
        \Lzero = \pi^2 \bdelta
    \end{equation}
    \tagP
\end{derivationbox}

\textbf{Interpretation:} $\pi^2 = \pi \times \pi$ reflects two independent
geometric constraints combining multiplicatively.

\subsection{Approach 3: Steiner Junction with Regularization}

\textbf{Physical picture:} The Y-junction (Steiner point) has three arms meeting
at 120° angles. The central hub is regularized at scale $\bdelta$.

\begin{derivationbox}[title=Steiner Junction Argument]
    The arm length $R$ (from hub to vertices) supports a standing wave mode.

    \textbf{Arm quantization:} For one full wavelength along the arm:
    \begin{equation}
        R = 2\pi \times \bdelta + \bdelta = (2\pi + 1)\bdelta \approx 7.3\bdelta
    \end{equation}

    \textbf{Including hub phase:} If the hub contributes phase $\pi$:
    \begin{equation}
        {L_0}_{\text{eff}} = 2\pi\bdelta + \pi\bdelta = 3\pi\bdelta \approx 9.4\bdelta
    \end{equation}

    \textbf{Observation:} $3\pi \approx 9.42$ is within 5\% of $\pi^2 \approx 9.87$.
    \tagP
\end{derivationbox}

\textbf{Status:} The Steiner argument gives $3\pi$, not exactly $\pi^2$. The
5\% difference may arise from:
\begin{itemize}
    \item Geometric corrections to the 120° ideal angles
    \item Finite size effects at the hub
    \item Higher-order terms in the 5D action
\end{itemize}

\subsection{Summary of Geometric Approaches}

\begin{center}
\begin{tabular}{llll}
\toprule
Approach & Result & Status & Notes \\
\midrule
Standing wave & $\pi^2 = 9.87$ & [P] & Assumes $R_5 = \pi\bdelta$ \\
Two-fold winding & $\pi^2 = 9.87$ & [P] & Assumes resonance $k_r = k_\theta$ \\
Steiner junction & $3\pi = 9.42$ & [P] & 5\% off from $\pi^2$ \\
\bottomrule
\end{tabular}
\end{center}

All approaches give $\Lzero/\bdelta \approx 9$--$10$, consistent with $\pi^2$.

% ============================================================================
\section{Potential-Theoretic Approach (5D)}
\label{sec:L0delta:potential}

An alternative derivation uses the 5D Green's function structure.

\subsection{Setup}

Consider the potential $\Phi$ sourced by the junction in 5D:
\begin{equation}
    \nabla^2_{5D} \Phi = -\rho_{\text{junction}}
\end{equation}

where $\rho_{\text{junction}}$ is the junction's ``topological charge'' density.

\subsection{Characteristic Scale}

The solution involves a length scale set by the geometry:
\begin{equation}
    \Phi(r, w) \sim \frac{1}{(r^2 + w^2)^{3/2}}
\end{equation}

where $r$ is the radial coordinate and $w$ is the fifth coordinate.

\textbf{Cutoffs:}
\begin{itemize}
    \item UV cutoff: $\bdelta$ (brane thickness)
    \item IR cutoff: $\Lzero$ (junction extent)
\end{itemize}

\subsection{Dimensional Argument}

The total ``charge'' enclosed by the junction:
\begin{equation}
    Q \sim \int_{\bdelta}^{\Lzero} \rho \, d^5x \sim \Lzero^5 / \bdelta^2
\end{equation}

For topological quantization ($Q = 1$ in appropriate units):
\begin{equation}
    \Lzero^5 \sim \bdelta^2 \quad \Rightarrow \quad \frac{\Lzero}{\bdelta} \sim \bdelta^{-3/5}
\end{equation}

\textbf{Status:} This approach does not directly yield $\pi^2$. It shows that
$\Lzero/\bdelta$ is determined by the 5D geometry, but the explicit value
requires solving the full boundary value problem.

\begin{edcopenbox}[title=Open: Rigorous 5D Derivation]
    \textbf{What is needed:}
    \begin{itemize}
        \item Solve the 5D Laplace/Poisson equation with junction boundary
              conditions
        \item Include brane tension term (Nambu-Goto contribution)
        \item Include bulk gravitational terms (GHY boundary conditions)
        \item Extract $\Lzero/\bdelta$ from the resulting field configuration
    \end{itemize}

    \textbf{Status:} Full calculation = [OPEN]. Available arguments give
    $\Lzero/\bdelta \sim O(10)$ but do not prove $\pi^2$ exactly.
    \tagOpen
\end{edcopenbox}

% ============================================================================
\section{Tension with Observables}
\label{sec:L0delta:tension}

\subsection{The $\pi^2$ vs $3\pi$ Question}

Different derivation approaches give slightly different values:
\begin{center}
\begin{tabular}{lll}
\toprule
Value & $\Lzero/\bdelta$ & $\SE/\hbar$ \\
\midrule
$\pi^2$ & 9.87 & 62.0 \\
$3\pi$ & 9.42 & 59.2 \\
Empirical (from $r_p$) & 9.33 & 58.6 \\
\bottomrule
\end{tabular}
\end{center}

\subsection{Impact on Metastable Lifetime}

The exponential sensitivity amplifies small differences:
\begin{equation}
    \frac{\tau(\pi^2)}{\tau(3\pi)} = \exp[(62.0 - 59.2)] = \exp(2.8) \approx 16
\end{equation}

A 5\% difference in $\Lzero/\bdelta$ leads to a factor of $\sim 16$ in $\taun$.

\subsection{Resolution}

The prefactor $A$ in $\taun = A \cdot (\hbar/\omegaz) \cdot \exp(\SE/\hbar)$
absorbs this ambiguity:
\begin{itemize}
    \item With $\Lzero/\bdelta = \pi^2$: need $A \approx 0.03$--$0.1$
    \item With $\Lzero/\bdelta = 9.33$: need $A \approx 0.8$--$1.0$
\end{itemize}

\textbf{Current approach:} We adopt $\pi^2$ as the geometric hypothesis [P] and
treat $A$ as a calibration parameter [Cal], checked against the requirement
that $A = O(1)$ for consistency.

% ============================================================================
\section{Epistemic Status}
\label{sec:L0delta:epistemic}

\begin{center}
\begin{tabular}{llll}
\toprule
\textbf{Claim} & \textbf{Status} & \textbf{Reference} & \textbf{Dependency} \\
\midrule
$\bdelta = \hbar/(2m_p c)$ & [I] & Eq.~\eqref{eq:delta_def} & Book I input \\
$\Lzero$ = junction extent & [Dc] & \S\ref{sec:L0delta:scales} & Definition \\
Standing wave gives $\pi^2$ & [P] & \S\ref{sec:L0delta:geometry} & Assumes $R_5 = \pi\bdelta$ \\
Two-fold winding gives $\pi^2$ & [P] & \S\ref{sec:L0delta:geometry} & Assumes resonance \\
Steiner gives $3\pi$ & [P] & \S\ref{sec:L0delta:geometry} & Hub phase assumption \\
$\Lzero/\bdelta = \pi^2$ & [P] & Eq.~\eqref{eq:L0delta_hypothesis} & Best motivated guess \\
\midrule
5D derivation of $\Lzero/\bdelta$ & [OPEN] & \S\ref{sec:L0delta:potential} & From full action \\
Proof that exactly $\pi^2$ & [OPEN] & — & vs.\ $3\pi$ or other \\
\bottomrule
\end{tabular}
\end{center}

% ============================================================================
\section{Falsifiability}
\label{sec:L0delta:falsify}

\begin{edcopenbox}[title=Falsifiability Hooks]
    \begin{enumerate}
        \item \textbf{Metastable lifetime mismatch:} If the full formula
              $\taun = A \cdot (\hbar/\omegaz) \cdot \exp[2\pi \times \pi^2]$
              cannot match $\taun \approx 880\second$ for $A \in [0.1, 10]$,
              the $\pi^2$ hypothesis is falsified.

        \item \textbf{Independent $\Lzero$ measurement:} If future experiments
              (e.g., anchor structure at high $Q^2$) measure $\Lzero$ directly
              and find $\Lzero/\bdelta \neq \pi^2 \pm 10\%$, the hypothesis fails.

        \item \textbf{5D action inconsistency:} If rigorous solution of the 5D
              boundary value problem gives $\Lzero/\bdelta$ significantly
              different from $\pi^2$, the geometric motivation is wrong.

        \item \textbf{$3\pi$ turns out correct:} If $3\pi = 9.42$ (from Steiner)
              is the true value rather than $\pi^2 = 9.87$, this changes
              $\taun$ predictions by factor $\sim 16$.
    \end{enumerate}
\end{edcopenbox}

% ============================================================================
% Observer-side projection
% ============================================================================

\begin{observerbox}
Observer-side projection note: quoted terms are measurement labels for the
3D/4D projection of the 5D objects defined in this chapter; they do not
imply any conventional mechanism.

\begin{itemize}
    \item \MetastableJunction{} $\leftrightarrow$ ``neutron'' (projection label)
\end{itemize}

What the observer records: the geometric ratio $\Lzero/\bdelta = \pi^2$
sets the magnitude of the Euclidean action, directly determining the
measured lifetime $\taun \approx 880\second$ at the projection level.
\end{observerbox}

% ============================================================================
\section{Summary}
\label{sec:L0delta:summary}

\begin{resultbox}
    The geometric ratio $\Lzero/\bdelta$ controls the instanton action:
    \begin{equation*}
        \frac{\SE}{\hbar} = 2\pi \times \frac{\Lzero}{\bdelta}
    \end{equation*}

    \textbf{Hypothesis [P]:} $\Lzero/\bdelta = \pi^2 \approx 9.87$

    \textbf{Geometric motivations:}
    \begin{itemize}
        \item Standing wave: $\Lzero = \pi R_5$, $R_5 = \pi\bdelta$ $\Rightarrow$ $\pi^2$
        \item Two-fold winding: radial $\times$ angular $= \pi \times \pi$
        \item Steiner junction: gives $3\pi \approx \pi^2$ (5\% difference)
    \end{itemize}

    \textbf{Status:} Plausible hypothesis [P], not rigorous derivation [Der].
    Full derivation from 5D action remains [OPEN].
\end{resultbox}

\bigskip
\noindent
\textbf{Forward references:}
\begin{itemize}
    \item Chapter~\ref{ch:tau_n} combines $\kappa = 2\pi$ and $\Lzero/\bdelta = \pi^2$
          to compute $\taun \approx 880\second$
\end{itemize}

\bigskip
\noindent
\textbf{Derivation chain status:}
\begin{equation*}
    \frac{\SE}{\hbar} = \underbrace{2\pi}_{\kappa \text{ [Dc], Ch.~\ref{ch:kappa}}}
              \times \underbrace{\pi^2}_{\Lzero/\bdelta \text{ [P], this chapter}}
              = 2\pi^3 \approx 62
\end{equation*}

\begin{edcopenbox}[title=Open Problem: Rigorous Derivation]
    \textbf{Goal:} Derive $\Lzero/\bdelta = \pi^2$ from the 5D action including:
    \begin{itemize}
        \item Bulk gravitational terms
        \item Brane tension (Nambu-Goto)
        \item Gibbons-Hawking-York boundary term
        \item Israel junction conditions
    \end{itemize}

    \textbf{Status:} [OPEN]. Current status is geometric motivation [P].
    \tagOpen
\end{edcopenbox}

% ----------------------------------------------------------------------------
% BRIDGE TO NEXT CHAPTER
% ----------------------------------------------------------------------------
\paragraph{Bridge to Chapter~\ref{ch:tau_n}.}
With $\kappa = 2\pi$~[Dc] and $L_0/\delta = \pi^2$~[P] now established,
the Euclidean action is fully determined: $\SE/\hbar = 2\pi^3 \approx 62$.
Chapter~\ref{ch:tau_n} assembles these components into the final metastable
lifetime prediction $\taun \approx 880\second$, completing the instanton
sector of Book IV.

